\section{Experimental Evaluation}

The results of the experimental evaluation for the developed model will be presented along four aspects: overall classification performance, per-class analysis using confusion matrix, global feature importance using SHAP value and quantitative assessment of the counterfactual recommendation engine.

\subsection{Overall Classification Performance}

Overall performance of the models deployed along with the baseline counterparts on the unseen test dataset with 178,385 instances is reported in Table~\ref{tab:overall}.

\begin{table}[htbp]
\centering
\caption{Overall model performance comparison.}
\label{tab:overall}
\begin{tabular}{lccccc}
\hline
\textbf{Model} & \textbf{Accuracy} & \textbf{Weighted F1} & \textbf{Macro F1} & \textbf{AUC (macro)} & \textbf{AUC (weighted)} \\
\hline
Naive baseline & 53.31\% & --- & --- & --- & --- \\
Logistic Regression & 77.67\% & --- & --- & --- & --- \\
\textbf{XGBoost} & \textbf{92.78\%} & \textbf{0.93} & \textbf{0.84} & \textbf{0.9864} & \textbf{0.9857} \\
\textbf{LightGBM} & \textbf{92.98\%} & \textbf{0.93} & \textbf{0.85} & \textbf{0.9870} & \textbf{0.9859} \\
\hline
\end{tabular}
\end{table}

The accuracy of both the gradient boosting techniques used is above 92\% for the five classes' classification problem, with the LightGBM slightly outperforming XGBoost by 0.20\%. Both models exhibit macro One-versus-Rest AUC ROC value higher than 0.986, which suggests a high level of class separation amongst all five classes. Importantly, both models significantly outperform the baseline of naive majority-class prediction at 53.31\%, which predicts Stable for all cases and the baseline of Logistic Regression at 77.67\%.

The discrepancy between weighted F1 (0.93) and macro F1 (0.84--0.85) measures arises from different levels of accuracy in classes of varying sizes --- an aspect further analyzed in Section~4.2. In this case, the macro F1 offers the more conservative measure of performance, as its calculations give equal weight to all classes irrespective of their frequency, and is therefore the measure to use when analyzing a system in which correct classifications of classes that occur less frequently (i.e., At-Risk and Thriving) have practical significance equivalent to their clinical significance.

With regard to comparative metrics, Altman et al.~\cite{4} cite 75\% accuracy in the Z-Score model in international environments, while Barboza et al.~\cite{5} show that machine learning models provide around 10 percentage point improvements compared to traditional models in binary bankruptcy prediction, which Wang et al.~\cite{14} corroborate. The lack of a precedent for classification models in five-class categories in the SBA dataset means numerical comparison is impossible; however, this finding itself validates the innovation in the taxonomy proposed. That said, the fact that the current study managed to obtain 92.78\%--92.98\% accuracy when faced with such a more complex five-class challenge indicates that the graduated viability taxonomy does not detract from predictive performance when compared to the binary taxonomy.

\subsection{Per-Class Report}

Classification reports per class for the XGBoost and LightGBM models, respectively, are found in Tables~\ref{tab:xgb_class} and~\ref{tab:lgbm_class}.

\begin{table}[htbp]
\centering
\caption{XGBoost per-class classification report (test set $n = 178{,}385$).}
\label{tab:xgb_class}
\begin{tabular}{lccccc}
\hline
\textbf{Class} & \textbf{Precision} & \textbf{Recall} & \textbf{F1-Score} & \textbf{Support} & \textbf{\% of Test Set} \\
\hline
Critical (0) & 0.84 & 0.80 & 0.82 & 28,909 & 16.21\% \\
At-Risk (1) & 0.69 & 0.40 & 0.51 & 2,197 & 1.23\% \\
Stable (2) & 0.94 & 0.95 & 0.95 & 95,100 & 53.31\% \\
Growing (3) & 0.94 & 0.98 & 0.96 & 31,070 & 17.42\% \\
Thriving (4) & 0.96 & 0.99 & 0.97 & 21,109 & 11.83\% \\
\hline
\end{tabular}
\end{table}

\begin{table}[htbp]
\centering
\caption{LightGBM per-class classification report (test set $n = 178{,}385$).}
\label{tab:lgbm_class}
\begin{tabular}{lccccc}
\hline
\textbf{Class} & \textbf{Precision} & \textbf{Recall} & \textbf{F1-Score} & \textbf{Support} & \textbf{\% of Test Set} \\
\hline
Critical (0) & 0.85 & 0.81 & 0.83 & 28,909 & 16.21\% \\
At-Risk (1) & 0.68 & 0.43 & 0.53 & 2,197 & 1.23\% \\
Stable (2) & 0.95 & 0.95 & 0.95 & 95,100 & 53.31\% \\
Growing (3) & 0.94 & 0.97 & 0.96 & 31,070 & 17.42\% \\
Thriving (4) & 0.96 & 0.99 & 0.97 & 21,109 & 11.83\% \\
\hline
\end{tabular}
\end{table}

Confusion matrices $5 \times 5$ are displayed in Tables~\ref{tab:cm_xgb} and~\ref{tab:cm_lgbm}. \textit{Note:} The confusion matrices provided below have been produced using a replicated data preprocessing pipeline ($n = 177{,}988$). Differences in class support between Tables~\ref{tab:xgb_class}--\ref{tab:lgbm_class} ($n = 178{,}385$) are caused by minor differences in data cleaning; they do not change qualitative misclassification patterns.

\begin{table}[htbp]
\centering
\caption{XGBoost confusion matrix (rows = actual, columns = predicted).}
\label{tab:cm_xgb}
\begin{tabular}{lccccc}
\hline
& \textbf{Critical} & \textbf{At-Risk} & \textbf{Stable} & \textbf{Growing} & \textbf{Thriving} \\
\hline
\textbf{Critical} & \textbf{23,439} & 7 & 4,424 & 668 & 0 \\
\textbf{At-Risk} & 10 & \textbf{997} & 2 & 845 & 429 \\
\textbf{Stable} & 5,179 & 2 & \textbf{89,426} & 475 & 329 \\
\textbf{Growing} & 213 & 282 & 58 & \textbf{30,089} & 783 \\
\textbf{Thriving} & 0 & 42 & 25 & 37 & \textbf{20,227} \\
\hline
\end{tabular}
\end{table}

\begin{table}[htbp]
\centering
\caption{LightGBM confusion matrix (rows = actual, columns = predicted).}
\label{tab:cm_lgbm}
\begin{tabular}{lccccc}
\hline
& \textbf{Critical} & \textbf{At-Risk} & \textbf{Stable} & \textbf{Growing} & \textbf{Thriving} \\
\hline
\textbf{Critical} & \textbf{23,450} & 25 & 4,362 & 701 & 0 \\
\textbf{At-Risk} & 11 & \textbf{1,012} & 3 & 852 & 405 \\
\textbf{Stable} & 4,640 & 9 & \textbf{89,943} & 481 & 338 \\
\textbf{Growing} & 248 & 294 & 98 & \textbf{29,935} & 850 \\
\textbf{Thriving} & 0 & 76 & 36 & 54 & \textbf{20,165} \\
\hline
\end{tabular}
\end{table}

A few trends can be observed from the per-class results:

\textbf{Stable, Growing, and Thriving Classes (F1 $\geq$ 0.95)} together make up 82.56\% of the test data and show good classification results. The Stable majority (53.31\%) is the biggest majority and shows a well-balanced precision-recall curve (0.94--0.95). It must be stated that the Stable majority causes a high accuracy of a naive classifier which predicts everything to be Stable (53.31\%). However, we consider the macro-F1 of 0.84--0.85, which does not depend on the class frequency, as the performance indicator and see that the performance level is still very high. The Growing and Thriving classes show outstanding recall levels (0.97--0.99), which shows that the classifier can detect almost all the enterprises in higher viability classes.

\textbf{Critical class (F1 = 0.82--0.83)} shows good accuracy with an identifiable confusion matrix pattern. Out of 28,909 Critical examples, around 4,400 are incorrectly categorized into Stable (15.3\%). This confusion across the two classes is because companies with health scores close to 25 have features similar to those on the lower boundary of the Stable class. It is noteworthy that the reverse mistake (stable-classified companies as Critical) is also seen (5,179 examples). In this case, we can conclude that the main confusion area lies in the Critical--Stable boundary. With respect to real-world applicability, misclassification from Critical to Stable means more expensive error, because it leads to lack of measures for struggling companies.

\textbf{At-Risk Class (F1 = 0.51--0.53)} is the class with the poorest performance. The reason behind this lies in two factors that work together. Firstly, the At-Risk class is just 1.23\% of the testing data (2,197 records), meaning that it is an extremely rare class, even though SMOTE was applied during the training process~\cite{12}. Secondly, the health score range that forms the At-Risk class (25--40) covers only 15 scores. These records lie at the boundary of feature distribution of both the neighboring classes. This is evident from the confusion matrix, where out of 2,197 instances classified into At-Risk, 845--852 are misclassified as Growing, while 405--429 are identified as Thriving. These records have not been misclassified in neighboring classes but in far-away classes, which confirms the difficulty of the At-Risk boundary region in terms of its ambiguity in feature space.

\subsection{SHAP Feature Attribution Analysis}

Global explainability of the model is estimated using average absolute SHAP values~\cite{17},~\cite{18} on a stratified random sample of 1,000 test cases from all five classes.

\begin{table}[htbp]
\centering
\caption{Global feature importance (average $|\text{SHAP}|$ on 1,000 samples, all classes).}
\label{tab:shap_global}
\begin{tabular}{clcc}
\hline
\textbf{Rank} & \textbf{Feature} & \textbf{Mean $|\text{SHAP}|$} & \textbf{Category} \\
\hline
1 & GrAppv & 1.5191 & Loan Structure \\
2 & Term & 1.4523 & Loan Structure \\
3 & RetainedJob & 0.6152 & Employment Impact \\
4 & SBA\_Appv & 0.4927 & Loan Structure \\
5 & CreateJob & 0.3778 & Employment Impact \\
6 & DisbursementGross & 0.1547 & Loan Structure \\
7 & NoEmp & 0.0611 & Enterprise Characteristics \\
8 & UrbanRural & 0.0471 & Enterprise Characteristics \\
9 & RevLineCr & 0.0399 & Loan Program \\
10 & NewExist & 0.0219 & Enterprise Characteristics \\
11 & LowDoc & 0.0124 & Loan Program \\
\hline
\end{tabular}
\end{table}

The following insights can be drawn. Firstly, \textit{GrAppv} (gross approved amount) and \textit{Term} (loan duration) feature combination accounts for a much larger share of SHAP contribution compared to all other features together, with the mean absolute SHAP value of 1.52 and 1.45 for each variable correspondingly. The high SHAP value for Gross Approved Amount is justified by the inclusion of gross amount in the health score calculation formula in addition to its indirect impact as the institutional enterprise capacity proxy (greater amount indicates greater lender trust and expectations regarding their ability to service the loan). Similar considerations justify the high contribution for Term due to its dual impact on health scores and market signals.

Secondly, employment features combination (consisting of Retained Job variable value of 0.62 and Create Job value of 0.38) ranks third in the list of most significant feature groups. Consistently with the domain knowledge, retained employment features have relatively higher predictive importance compared to creating jobs.

Thirdly, the loan program characteristics, (\textit{RevLineCr} = 0.04, \textit{LowDoc} = 0.01), show very little contribution from SHAP, indicating that the loan program administrative attributes play no role in the classification of loan program viability.

In order to show the local application of global importance, Table~\ref{tab:shap_instance} shows the SHAP breakdown of a sample Critical class prediction.

\begin{table}[htbp]
\centering
\caption{Instance-level SHAP contributions for a Critical-class prediction.}
\label{tab:shap_instance}
\begin{tabular}{lcc}
\hline
\textbf{Feature} & \textbf{SHAP Value} & \textbf{Direction} \\
\hline
Term & +0.5844 & Risk-amplifying \\
DisbursementGross & +0.0549 & Risk-amplifying \\
LowDoc & +0.0124 & Risk-amplifying \\
SBA\_Appv & $-$0.0102 & Risk-mitigating \\
NoEmp & $-$0.0157 & Risk-mitigating \\
RetainedJob & $-$0.0454 & Risk-mitigating \\
\hline
\end{tabular}
\end{table}

In this case, the short-term lending period is the main factor behind the Critical classification (SHAP = +0.5844), reflecting its importance as the second-most important variable worldwide. The number of employees hired is partly offsetting the impact of the evaluation process (SHAP = $-$0.0454), which means that the model sees proof of the company's ability to operate.

\subsection{Evaluation of the Counterfactual Recommendation Engine}

For the quantitative evaluation of the counterfactual recommendation engine~\cite{20},~\cite{21}, the engine is applied to 200 test cases selected in a systematic manner: 100 Critical cases and 100 At-Risk cases. In each test case, the engine must find the minimum perturbation of features needed to reach the next highest viability class.

\begin{table}[htbp]
\centering
\caption{Counterfactual engine performance ($n = 200$).}
\label{tab:cf_eval}
\begin{tabular}{lcc}
\hline
\textbf{Metric} & \textbf{Critical $\rightarrow$ At-Risk} & \textbf{At-Risk $\rightarrow$ Stable} \\
\hline
Sample size & 100 & 100 \\
Feasibility rate & 100/100 (100\%) & 100/100 (100\%) \\
Phase 1 (single-feature) & 82/100 (82\%) & 41/100 (41\%) \\
Phase 2 (pairwise) & 0/100 & 0/100 \\
Already at/above target & 18/100 & 59/100 \\
Mean features changed & 1.00 & 1.00 \\
Most recommended feature & Term (100\%) & Term (100\%) \\
\hline
\end{tabular}
\end{table}

There are many highlights in this study. The first point to be emphasized is that the algorithm has attained \textbf{100\% feasibility} in both cases, which means that all the Critical as well as At-Risk items have at least one way to improve in terms of search space. This is a highly practical implication because there will always be at least one suggestion for improvement from the algorithm to the company.

Secondly, in regard to the Critical instances that need to be altered, \textbf{82\% are handled through single-feature changes} (Phase~1), whereas 18\% of the cases refer to those where the prediction of the model with the original (unmodified) features already satisfies the threshold criteria. That is, there is no need for a counterfactual case in such scenarios. The latter type of result is especially applicable to the At-Risk class, as 59\% of its cases do not require any alteration. The latter results can be attributed to the recall limitation of the At-Risk category (0.40--0.43) since a considerable number of At-Risk patients' data is classified into a more significant category (Growing or Thriving). No instances needed pairwise feature alterations (Phase~2), thus demonstrating the minimality characteristic of the recommendation engine.

Third, the \textbf{Term is the most recommended feature}, being present in 100\% of all recommendation generations. Such a result shows high internal consistency with the SHAP analysis (Section~4.3), where Term is the second-most important global feature. The coincidence of the identified risk factors using SHAP with counterfactual recommendations demonstrates that the explainability and prescriptive parts of the model generate consistent results.

Fourth, the recommendation engine works within feasible latency boundaries. The mean latency of the prediction process is \textbf{0.99~ms} (p95: 1.64~ms), while the mean latency of the counterfactual generation process is \textbf{2.03~ms} (p95: 4.77~ms).

\begin{table}[htbp]
\centering
\caption{Representative counterfactual: Critical $\rightarrow$ At-Risk.}
\label{tab:cf_example}
\begin{tabular}{lccc}
\hline
\textbf{Feature} & \textbf{Original Value} & \textbf{Recommended Value} & \textbf{Change} \\
\hline
Term (months) & 36 & 48 & $\uparrow$ Increase by 12 months \\
\hline
\end{tabular}
\end{table}

In the current example, the only adjustment needed to move from the Critical category to At-Risk is increasing the repayment period by 12 months. The proposed course of action is feasible in terms of loan restructuring and aligns with our understanding that increased repayment period will lighten debt service payment burdens.
